\documentclass[11pt]{article}
\usepackage[margin=1in]{geometry}
\usepackage[T1]{fontenc}
\usepackage[utf8]{inputenc}
\usepackage{lmodern}
\usepackage{hyperref}
\usepackage{xcolor}
\usepackage{longtable}
\usepackage{booktabs}
\usepackage{array}
\usepackage{parskip}

\title{Manuscript Revision Map}
\author{}
\date{March 15, 2026}

\begin{document}
\maketitle

\section*{Purpose}
This document records the most important wording changes made between the earlier manuscript
version and the current revised file \texttt{warpdrive\_clean\_submission.tex}. The goal is to make
the editorial logic explicit: what the earlier wording suggested, what the revised wording now says,
and why the change was necessary after auditing the regenerated evidence.

The changes summarized here are grounded in the cleaned optimization pipeline and in the associated
audit notes:
\begin{itemize}
\item regenerated optimized bundles for the single-shell and double-shell cases,
\item direct evaluations of the manuscript target parameters under the same cleaned pipeline,
\item the synthesis note \texttt{phase\_c\_synthesis.md}.
\end{itemize}

\section*{How to Read This Map}
Each entry below is organized into three parts:
\begin{itemize}
\item \textbf{Earlier emphasis}: what the pre-revision manuscript effectively claimed,
\item \textbf{Revised emphasis}: what the new manuscript now states,
\item \textbf{Reason for change}: why the earlier wording was too strong or too specific.
\end{itemize}

\renewcommand{\arraystretch}{1.25}
\begin{longtable}{p{0.17\textwidth} p{0.24\textwidth} p{0.24\textwidth} p{0.27\textwidth}}
\toprule
\textbf{Section} & \textbf{Earlier Emphasis} & \textbf{Revised Emphasis} & \textbf{Reason for Change} \\
\midrule
\endfirsthead
\toprule
\textbf{Section} & \textbf{Earlier Emphasis} & \textbf{Revised Emphasis} & \textbf{Reason for Change} \\
\midrule
\endhead
\bottomrule
\endfoot

Abstract &
The abstract leaned toward the idea that WEC/NEC were close to saturation and that the
SEC was effectively positive on the sampled slices. &
The abstract now states a hierarchy of residual constraints: $\rho$ remains non-negative by
construction, the DEC is the tightest residual condition, and the SEC is easier than the
DEC only in aggregate. &
The regenerated bundles and the manuscript-target evaluations both show that WEC/NEC remain
substantially constrained, while the SEC is not uniformly positive on the sampled slices. \\[0.2cm]

Introduction &
The introduction risked sounding like the optimization strongly rehabilitated the classical
energy conditions. &
The introduction now presents the method as a \emph{physics-penalized parametric optimization}
over an analytic ansatz, with no claim of full energy-condition recovery. &
This phrasing matches what the code actually does and avoids overselling the numerical workflow as
either a PINN or a proof of global energy-condition satisfaction. \\[0.2cm]

Single-shell discussion &
The earlier single-shell discussion suggested that WEC/NEC become nearly saturated and that the
remaining obstruction is mostly DEC-only. &
The revised text says that DEC remains the strongest residual obstruction, but WEC/NEC are
still substantially constrained and the SEC is mixed on the sampled slices. &
Both the optimized single-shell bundle and the direct evaluation of the manuscript parameters show
only moderate WEC/NEC success fractions, not a near-saturation regime. \\[0.2cm]

Double-shell discussion &
The earlier double-shell discussion emphasized localized deficits near the inner interface but
still gave the impression that the condition hierarchy was visually cleaner than it really is. &
The revised discussion keeps the qualitative localization story while making clear that the
double-shell case remains more constrained and that the SEC is not uniformly positive. &
The regenerated double-shell bundle still supports the geometric interpretation, but the cleaned
diagnostics do not justify claiming that the softer conditions are visibly or numerically close to
full satisfaction. \\[0.2cm]

SEC slice captions &
The earlier captions for the meridional SEC figures read as though positivity on the shown slices
had effectively been established. &
The revised captions say that the SEC is comparatively milder than the DEC in aggregate but
can still show mixed-sign structure on sampled slices. &
This was one of the most important corrections because the regenerated evidence does not support a
uniformly positive slice interpretation. \\[0.2cm]

Optimization curves &
The earlier captions and discussion described a strong drop in the physics loss followed by a clean
plateau and an initial transient in the parameters. &
The revised text describes the visible main run more carefully: pretraining places the optimization
near a stable basin, the late-time total loss is dominated by the inverse-alpha term, and the
main parameters remain in narrow bands. &
This is closer to what the regenerated runs show. The earlier wording read too much dynamical
story into curves whose visible main phase is already close to a plateau. \\[0.2cm]

Success-fraction captions &
The earlier wording implied a large numerical gap between NEC/WEC and DEC. &
The revised wording avoids calling the fractions ``high'' and instead states that the curves settle
into topology-dependent plateaus, with DEC remaining the tightest residual test. &
The cleaned runs still support DEC as the dominant obstruction, but not the stronger claim that
NEC/WEC enter an obviously high-success regime. \\[0.2cm]

Conclusion &
The earlier conclusion could be read as a stronger endorsement of the standard energy conditions
than the new evidence warrants. &
The revised conclusion focuses on the limited but defensible result: the surviving deficits are
stress-dominated, localized, and structured by topology, with DEC as the hardest residual test. &
This lets the paper keep its main constructive result while staying aligned with the regenerated
bundles and the direct paper-parameter evaluations. \\

\end{longtable}

\section*{Preserved Claims}
The revision does not overturn the core analytic structure of the paper. The following claims were
preserved and are still supported by the cleaned workflow:
\begin{itemize}
\item the construction should be described as a constrained or penalized optimization over a
low-dimensional analytic ansatz,
\item the seed is organized so that $\rho \ge 0$ is retained under the restricted uniform,
time-independent translations considered here,
\item the dominant residual obstruction is stress-sector dominated, with the DEC typically the
tightest of the sampled conditions,
\item topology matters: the double-shell case is more constrained and redistributes curvature more
subtly than the single-shell case.
\end{itemize}

\section*{Softened Claims}
The following claims were explicitly weakened because the regenerated evidence does not support the
stronger wording:
\begin{itemize}
\item ``high'' WEC/NEC success fractions,
\item uniformly non-negative WEC$_{\min}$ on the sampled slices,
\item uniformly positive SEC throughout the sampled shell regions,
\item a strong visible transient in the main optimization run,
\item a narrative in which the visible curves are dominated by a large drop in physics loss.
\end{itemize}

\section*{Editorial Principle}
The revised manuscript follows a simple principle: when the cleaned pipeline and the direct
manuscript-parameter evaluations support a weaker claim than the older wording suggested, the paper
now adopts the weaker but reproducible statement.

\end{document}
